% 小芽项目报告 —— 排版设置
% 编译：tectonic（XeTeX 引擎）+ ctex
%
% 注意：
% 1. 不设 \XeTeXlinebreaklocale / \XeTeXlinebreakskip。断行由 xeCJK 的 \CJKglue 处理，
%    手设这两个裸 primitive 在 ctex 下是不必要的（源码里除 CJLineBreak 外并无引用）。
% 2. 不用 titlesec。加载它会完全接管章节标题，ctex 自带的标题间距设置会失效。
%    章节样式改用 ctex 自己的 \ctexset。
% 3. 不用 \sloppy。它的作用是放宽断行阈值，代价是把不良断行的警告一起盖掉，
%    等于把问题藏起来。真正该做的是让长串可断（\url、表格用 p{} 列）并把内容写短。
% 4. 不写 \setlength\parindent。ctex 默认 autoindent=true，已经等于 2\ccwd，
%    而 2\ccwd 包含字距，与 2em 并不等价，手写反而容易错。
\documentclass[UTF8, scheme=chinese, a4paper, zihao=-4]{ctexart}

\usepackage{geometry}
\geometry{a4paper, top=2.5cm, bottom=2.5cm, left=2.6cm, right=2.6cm}

\usepackage{graphicx}
\usepackage{booktabs}
\usepackage{array}
\usepackage{enumitem}
\usepackage{xurl}              % 允许 URL 任意位置断行
\usepackage[hidelinks]{hyperref}

\linespread{1.3}               % ctex 的 scheme=chinese 默认行距

% 章节样式：用 ctex 自己的接口，不引 titlesec
\ctexset{
  section/format    = \normalfont\large\bfseries,
  subsection/format = \normalfont\normalsize\bfseries,
  section/beforeskip  = 1.4em,
  section/afterskip   = 0.6em,
  subsection/beforeskip = 1em,
  subsection/afterskip  = 0.4em,
}

\pagestyle{plain}

\newcommand{\code}[1]{\texttt{\small #1}}
\newcommand{\file}[1]{\texttt{\small #1}}

% 定宽表格列：内容短，不设 \sloppy 也不会溢出
\newcolumntype{L}[1]{>{\raggedright\arraybackslash}p{#1}}
